%------------ Debut Packages ----------------

\usepackage[T1]{fontenc} 
% Permet une meilleure gestion des accents et caractères spéciaux à l'impression (encodage des fontes en T1)

\usepackage[utf8]{inputenc} 
% Définit l'encodage des fichiers source en UTF-8 (utile pour les caractères accentués)

\usepackage{lmodern} 
% Utilise la police Latin Modern (amélioration de la police Computer Modern par défaut)

\usepackage[french]{babel} 
% Adapte le document au français (traduction automatique de mots comme "Chapter" → "Chapitre", gestion des césures...)

\usepackage{xcolor} 
% Permet d'utiliser les couleurs dans le document (texte, fond, tableaux, etc.)

\usepackage{multirow} 
% Pour fusionner plusieurs lignes dans un tableau

\usepackage{makecell} 
% Pour créer des cellules de tableaux personnalisées (ex : texte sur plusieurs lignes dans une cellule)

\usepackage{float} 
% Permet un placement plus précis des figures et tableaux (avec [H] pour figer à l’endroit exact)

\usepackage{tablefootnote} 
% Permet d’ajouter des notes de bas de page dans les tableaux

\usepackage{hyperref} 
% Pour insérer des liens cliquables (URL, références internes, etc.)

\usepackage{bookmark} 
% Améliore la gestion des signets PDF (utile en complément de hyperref)

\usepackage{pdfpages} 
% Permet d'inclure des pages PDF externes dans ton document

\usepackage[linesnumbered,ruled,vlined]{algorithm2e} 
% Pour écrire des algorithmes formatés (avec numérotation des lignes, règles et indentations)

\usepackage{lipsum} 
% Pour générer du faux texte (texte de remplissage, utile en phase de test ou maquette)

\usepackage{listings} 
% Pour insérer du code source avec mise en forme (coloration syntaxique, encadré, etc.)

\usepackage{subcaption} 
% Pour insérer plusieurs figures ou tableaux côte à côte avec des sous-titres (subfigures/subtables)

\usepackage{lscape} 
% Permet de faire pivoter certaines pages en mode paysage (utile pour les grands tableaux ou schémas)

\usepackage{url} 
% Pour gérer l’affichage des URL (utile pour que les adresses web soient bien formatées)

%------------ Fin des packages classiques ----------------

%------------ Entêtes et pieds de page ----------------

\usepackage{fancyhdr} 
% Permet de personnaliser les entêtes et pieds de pages

\pagestyle{fancy} 
% Active le style fancy pour toutes les pages standards du document

\renewcommand{\chaptermark}[1]{\markboth{\chaptername \ \thechapter.\ #1}{}} 
% Personnalise l'entête de gauche (nom du chapitre) : "Chapitre 1. Titre", sans tout mettre en majuscule

\renewcommand{\sectionmark}[1]{\markright{\thesection.\ #1}} 
% Personnalise l'entête de droite (nom de la section courante) : "1.1 Titre de section", sans majuscules

\renewcommand{\thesection}{\arabic{section}} 
% Modifie le style de numérotation des sections (exclut le numéro de chapitre, utile dans les rapports courts)

%------------ Fin Packages ----------------